% ch14: 原文 GLDraft_D.tex 行 3426–3539；附录 A
% 只译散文（标题、正文、图题、脚注）；公式、\label、\ref、\eqref、\cite、图文件名保持不变。

\appendix


\section{沿典范余迷向膜的局部算符与 $\varphi {\cal F}$ 理论中的量子修正}\label{coho}
本附录讨论变形 Neumann 边界条件（deformed Neumann boundary condition）的全部边界局部算符空间，
该边界条件对应于典范余迷向膜（canonical coisotropic brane）${\mathcal B}_{\mathrm{cc}}$。这个问题最初在 \cite{KW} 第 12.4 节中处理过。
这里我们将推广并部分修正那里的分析。

在经典层面，相关边界局部算符空间的 BRST 上同调如下。一类生成元是局部算符 ${\mathcal P}(\phi_z)$，
其中 ${\mathcal P}$ 是规范群 $G$ 的李代数 $\mathfrak g$ 上的不变多项式。第二类生成元是由
${\mathcal P}'(\phi_z+\epsilon\psi_z)={\mathcal P}(\phi_z)+\epsilon {\mathcal P}'(\phi_z,\psi_z)$ 定义的算子
${\mathcal P}'(\phi_z,\psi_z)$，其中 $\epsilon$ 是一个奇参数。
结果表明，经典地，${\mathcal B}_{\mathrm{cc}}$ 的边界局部算符的 BRST 上同调由算子 ${\mathcal P}$、
${\mathcal P}'$ 及其对 $z$ 的导数生成。这一事实稍后就会解释。
人们预期，完整量子理论中的 BRST 上同调是作用在经典上同调上的某个进一步的微分 $Q$ 的上同调。这里的目标
是证明这个微分消失，从而量子理论中的上同调与经典情形相同。如第 \ref{hh} 节所解释，这尤其蕴涵
交换的 Hitchin 哈密顿量可以量子化为交换的全纯微分算子。

由于上同调经典地由形如 $\mathcal P$ 与 $\mathcal P'$ 的算子及其对 $z$ 的导数生成，
又由于 $\partial_z$ 与 $Q$ 交换，只需证明对 $\mathfrak g$ 上任一不变多项式 $\mathcal P$ 都有 $[Q,\mathcal P]=\{Q,\mathcal P'\}=0$。
为此，可以寻找一个能强制这一结果的对称性。

在 \cite{KW} 第 12.4 节中已经注意到，$Q$ 与一个时间反演对称性交换，后者作用为
\begin{equation}\label{tadd}
{\sf T}: t \to - t \qquad \qquad \varphi_z \to - \varphi_z \qquad \qquad \psi_z \to \psi_z.
\end{equation}
对 $G=\mathrm{SU}(2)$，这个对称性就足以推出我们想要的结果。对这个群，不变多项式环由 $\mathcal P(\phi_z)={\mathrm{Tr}}\,
\phi_z^2$ 生成，所以只需对这个特定的 $\mathcal P$ 证明 $[Q,\mathcal P]=\{Q,\mathcal P'\}=0$。为证明 $[Q,\mathcal P(\phi_z)]=0$，我们注意到，
上同调中与 $[Q,\mathcal P(\phi_z)]$ 具有相同自旋和费米子数\footnote{我们用到了如下事实：$\phi_z,\psi_z$ 与 $\partial_z$ 的自旋都是 1，而
$\phi_z$ 的费米子数为 0、$\psi_z$ 的费米子数为 1。$Q$ 的自旋为零、费米子数为 1。}的算子只有 $\mathcal P'(\phi_z,\psi_z)$，
但 $\mathcal P(\phi_z)$ 在 $\sf T$ 下为偶，而 ${\mathcal P}'(\phi_z,\psi
_z)$ 为奇。
同样成立的是 $\{Q,\mathcal P'\}=0$，因为经典
上同调中不含任何与 $\{Q,\mathcal P'\}$ 具有相同自旋和费米子数的非零算子。

这一论证对秩更高的群一般并不充分。一旦考虑到本问题中的经典上同调是由 $z$ 导数 $\partial_z$ 以及算子 $\phi_z$ 与 $\psi_z$ 一同构造的，
就会发现确实存在一些候选算子，它们可以出现在
公式 $[Q,\mathcal P]=?$ 与 $\{Q,\mathcal P'\}=?$ 的右端，并且与 $\sf T$ 对称性相容。然而，还有另一个相关的对称性。这就是对称性 $\sf C$：
物理上称为电荷共轭（charge conjugation），数学上称为 Chevalley 对合（Chevalley involution）。任何单李代数 $\mathfrak g$ 的根图在
把每个根向量变为 $-1$ 倍的操作 $\sf C$ 下不变。由于 $\mathfrak g$
可由根向量的长度与它们之间的夹角重构，
$\sf C$ 延拓为 $\mathfrak g$ 的对合，从而是相应李群 $G$ 的对合。这个对称性就是电荷共轭或 Chevalley 对合。
（对某些群，例如 ${\mathrm{SO}}(2n+1)$，它是内自同构；对其他群，特别是 ${\mathrm{SU}}(n),\,n>2$、
${\mathrm{SO}}(2n)$ 与 $E_6$，它是外自同构。对 $\sf C$ 为内的 $G$，仅用 $\sf T$ 就能完成下面的论证。）\zy{译者按（无外部依据）：$\sf C$ 即 Chevalley 对合，在 Chevalley 基下可写成 $X_\alpha\mapsto -X_{-\alpha}$、$h\mapsto -h$，故它在极大环面（Cartan 子代数）上恰为 $-1$。原文“把每个根向量变为 $-1$ 倍”应理解为根图上的反射 $\alpha\mapsto -\alpha$；若照字面把每个根向量都乘以 $-1$，一般并不给出李代数自同构。}

由于 $\mathfrak g$ 上的不变多项式 $\mathcal P$ 在极大环面的李代数上有非平凡的限制，
可以通过限制到极大环面来判断 $\mathcal P(\phi_z)$ 在 $\sf C$ 下是偶还是奇。结果人们发现，
在所有情形下都有 $\sf C \mathcal P(\phi(z)) \sf C^{-1}=\mathcal P(-\phi_z)$（尽管 $\sf C \phi_z\sf C^{-1}=-\phi_z$ 并不成立）。\zy{译者按（无外部依据）：此处系原文校勘。原文作 “one one finds”，重复了一个 “one”；同句又作 $\mathcal P(\phi(z))$，按上下文应为 $\mathcal P(\phi_z)$。此外式 (\ref{tadd}) 用 $\varphi_z$，正文其余地方多用 $\phi_z$，属原稿记号不统一。译文对公式一律照录，未作改动。}
同样，
$\sf C\mathcal P'(\phi_z,\psi_z)\sf C^{-1} =\mathcal P'(-\phi_z,-\psi_z)=-\mathcal P'(-\phi_z,\psi_z)$。把它与式 (\ref{tadd}) 结合，
可知在所有情形下
\begin{equation}\label{wadd} \sf C\sf T \mathcal P(\phi_z)(\sf C\sf T)^{-1}=\mathcal P(\phi_z),~~~
  \sf C\sf T \mathcal P'(\phi_z,\psi_z)(\sf C\sf T)^{-1}=-\mathcal P'(\phi_z,\psi_z). \end{equation}
因此，用算子 $\mathcal P$、$\mathcal P'$ 及其 $z$ 导数中的多项式来构造经典上同调，我们得知 $\sf C\sf T$
在经典上同调中费米子数为 $k$ 的部分（即对算子 $\mathcal P'$ 及其
导数而言为 $k$ 次齐次的部分）上作用为 $(-1)^k$。由于 BRST 微分 $Q$ 与 $\sf C\sf T$ 交换，可知该算子必然平凡地作用在本问题的经典
上同调上。的确，公式 $[Q,\mathcal P]=?$ 或 $\{Q,\mathcal P'\}=?$ 的左端与右端
在 $\sf C\sf T$ 下会以相反的方式变换。

\begin{yizhu}[译者按（无外部依据）：$\sf C\sf T$ 为何排除量子修正]
原文只用一句话带过“$Q$ 与 $\sf C\sf T$ 交换”；从这里到结论“$Q$ 在经典上同调上为零”还差一步分次线性论证，这里补出。设 $Q$ 是作用在经典上同调上的进一步微分，次数为 1（费米子数为 1），且满足 $Q\,\sf C\sf T=\sf C\sf T\,Q$。由式 (\ref{wadd})，$\sf C\sf T$ 在费米子数为 $k$ 的齐次部分上等于 $(-1)^k$。若 $x$ 为 $k$ 次齐次元而 $Qx=y\ne 0$，则 $y$ 为 $k+1$ 次齐次元；由 $Q(\sf C\sf T x)=\sf C\sf T(Qx)$ 得 $(-1)^k y=(-1)^{k+1}y$，故 $y=0$，矛盾。于是 $Q\equiv 0$，量子上同调与经典上同调一致。这步隐含两条输入：其一，$\sf C\sf T$ 在经典上同调上的作用由式 (\ref{wadd}) 确定；其二，$\sf C\sf T$ 是量子理论的对称性，因而与量子 BRST 微分交换；若缺后者，“交换”只是假设，论证不成立。对 $\{Q,\mathcal P'\}=?$ 同理：两端在 $\sf C\sf T$ 下符号相反，故右端只能为零。
\end{yizhu}

在本附录余下的部分，我们将解释如下断言：在经典层面，相关局部算符空间的 BRST 上同调
由算子 $\mathcal P$、$\mathcal P'$ 及其导数生成。
我们在（解析延拓的）三维全纯-拓扑 $\varphi {\cal F}$ 理论的 BV 形式中研究这一问题。
方便的做法是采用带有奇坐标 ${\mathrm d} t$ 与 ${\mathrm d} \bar z$ 的超空间，这样超场就成为各阶形式之和。
物理场、鬼场与反场可以汇集到两个超场中，我们记之为 ${\cal B}$ 与 ${\cal A}$，其作用量为
\begin{equation}
\int \mathrm{Tr} \,{\cal B} \,{\cal F}_{\cal A}
\end{equation}
这与三维 Chern–Simons 理论的 BV 形式完全类似：那里对一个由各阶形式之和构成的超场
采用标准的 Chern–Simons 作用量。

${\cal B}$ 的 0-形式分量是 $\varphi_z$。${\cal A}$ 的 0-形式分量是鬼场 $c$，而 1-形式（只有 ${\mathrm d}t$ 与 ${\mathrm d} \bar z$ 分量）是规范联络。鬼场 $c$ 本身不允许出现在局部算符的构造中：只允许它的导数出现。这是紧致规范群的规范场论中 BRST 形式的一个标准特征。
全纯 $\varphi {\mathcal {F}}$ 理论中的全纯导数 $\partial_z c$ 扮演完整四维规范场论中 $\psi_z$ 的角色。

经典地，在自由理论中，BRST 微分来自动能项，并涉及偏微分
\begin{equation}
\mathrm{d} \equiv {\mathrm d}t \frac{\partial}{\partial t}+ {\mathrm d} \bar z \frac{\partial}{\partial \bar z}.
\end{equation}
两个超场都被 $Q_{\mathrm{free}} + \mathrm{d}$ 消灭。自由理论的上同调由
$\varphi_z$、$\psi_z$ 及其全纯导数或 $\partial_z$ 导数的 $G$-不变多项式组成。\footnote{动能项把超场的 $0$-形式分量
与 $1$-形式分量配对。这意味着 $0$-形式分量中的这类多项式有良好定义，没有重整化歧义。这也意味着量子自由理论的 BV Laplacian 不影响最终答案。}

经典 $\varphi \mathcal F$ 理论中的上同调与自由理论的上同调并不相同，因为
相互作用在树图层面就已经改变 BRST 微分。基本超场的树图层面 BRST 变分为
\begin{equation}
\left[Q_{\mathrm{tree}} + \mathrm{d}_{\cal A} \right] {\cal B} = 0 \qquad \qquad Q_{\mathrm{tree}} {\cal A} + {\cal F}_{\cal A} = 0
\end{equation}

经典（或树图层面）上同调可以这样计算：先过渡到自由理论的上同调，再把微分的非线性部分
限制到自由上同调上。这是分析扭化超对称场论的标准步骤。
取自由理论的上同调，相当于
丢弃 $D_t$ 与 $D_{\bar z}$ 导数以及场的高阶形式分量（只留下 $\partial_z$ 导数）。
所得的微分形如
\begin{equation}
Q \varphi_z = [c,\varphi_z] \quad \qquad Q c = \frac12 [c,c]
\end{equation}

这个上同调问题可以数学地表述为相对李代数上同调问题 $H^*(\mathfrak{g}[z],\mathfrak{g};S \mathfrak{g}^*[z])$。
底层的李代数 $\mathfrak{g}[z]$ 与鬼场 $c$ 及其全纯导数相联系。问题相对于 $\mathfrak{g}$，是因为
我们略去了 $c$ 场本身。上同调取值于 $S \mathfrak{g}^*[z]$，它表示 $\varphi_z$ 及其导数的多项式。
\zy{译者按（无外部依据）：相对李代数上同调 $H^*(\mathfrak g[z],\mathfrak g;S\mathfrak g^*[z])$ 指在 $\mathfrak g[z]$ 的 Chevalley–Eilenberg 复形中只取对 $\mathfrak g$ 不变、且在 $\mathfrak g\subset\mathfrak g[z]$ 上基本（缩并为零）的余链，即 $\mathrm{Hom}_{\mathfrak g}\bigl(\wedge^q(\mathfrak g[z]/\mathfrak g),S\mathfrak g^*[z]\bigr)$ 的上同调；原文所说的“略去 $c$ 场本身”，正对应这里的“相对”。}

这个上同调问题由 \cite{FGT} 中的定理 B 解决，结论正是前面所断言者：经典上同调由算子
$\mathcal P$、$\mathcal P'$ 及其 $z$ 导数生成。正如我们已解释的，存在一个对称性 $\sf {CT}$，它保证这一结果没有量子
修正。


我们所描述的 $A$-膜 $\mathcal B_{\mathrm{cc}}$ 的边界局部算符空间，
与 $S$-对偶的 oper 边界条件处的局部算符空间一致。
的确，oper 膜的 $B$-模型自同态是 oper 流形上的全纯形式，而 oper 流形是一个仿射向量空间，由
联系于 $\mathfrak g$ 上不变多项式 $\mathcal P$ 的算子参数化。对 oper 流形上的每个仿射坐标 $w$（它联系于某个
$\mathcal P$），都有相应的 1-形式 ${\mathrm d} w$，以同样的方式联系于 $\mathcal P'$。于是 oper 流形上的全纯形式以
自然的方式对应于 $A$-模型中所得到的东西。
